\documentclass{strategynet-report}

\hypersetup{
  pdftitle={Volt: Portfolio Control},
  pdfauthor={research@strategynet.ai}
}

\newcommand{\VOLT}{\textsf{Volt}}
\newcommand{\COG}{\texttt{cog}}
\newcommand{\E}{\mathbb{E}}

\tikzset{
  controlbox/.style={
    draw=SNLine,
    fill=SNPaper,
    rounded corners=2pt,
    text width=32mm,
    minimum height=21mm,
    align=center,
    inner sep=2.6mm,
    font=\sffamily\fontsize{8}{9.5}\selectfont
  },
  controllearn/.style={
    controlbox,
    draw=SNSignalDark,
    fill=SNSignal!6
  },
  controlflow/.style={
    -{Latex[length=2mm]},
    draw=SNInk,
    line width=0.75pt
  },
  controlfeedback/.style={
    -{Latex[length=2mm]},
    draw=SNSignalDark,
    dashed,
    line width=0.75pt
  }
}

\title{\VOLT{}: Portfolio Control}
\SNSubtitle{Managing options as state, risk, and available actions change}
\author{\href{mailto:research@strategynet.ai}{research@strategynet.ai}}
\SNDocumentType{Research architecture}
\SNVersion{Version 0.2}
\SNDate{July 2026}
\SNStatus{Open-source working paper}
\SNShortTitle{Volt: Portfolio Control}

\begin{document}
\maketitle

\begin{abstract}
Opening an options strategy is only the first decision.  Each hedge, close,
roll, exercise, assignment, and market move changes the position, its risk,
and the actions that remain available.  A policy trained on a static payoff
table misses that path dependence; a policy allowed to redefine contract or
accounting rules cannot be trusted.

This paper proposes a planning and reinforcement-learning layer that consumes
a graph of the current portfolio, its relationships, candidate actions, and
observed state.  \VOLT{} exposes that structure to the learner while remaining
responsible for contract meaning, fills, settlement, transitions, and replay.
The public crate already provides the deterministic runtime, scenario paths,
journals, and a static portfolio graph needed to begin such an environment.
It does not yet provide account policy, a reward model, a trained policy, or a
rollout scheduler.  A future \COG{} compiler could construct and coordinate
the graph, candidate, and legal-action artifacts efficiently.  QUBO is
considered only as a possible rollout-budget allocator and would need to beat
simpler schedulers under the same budget.
\end{abstract}

\SNtableofcontents

\part{Control Problem}

\section{Portfolio state}

Static discovery asks which admissible portfolio looks attractive at one
decision time.  Portfolio control asks what to do next, observes the
consequences, and asks again.  Closing a short option may free capital but
realize the spread; a hedge alters delta and financing; exercise consumes a
right and may create cash or underlying inventory; assignment changes the
account without being chosen by its owner.  A useful state must carry these
consequences forward rather than recomputing the position from a terminal
payoff diagram.

In the proposed control interface, \COG{} would emit
\[
  s_t=(G_t,A_t,M_t,p_t),
\]
where \(G_t\) describes the current portfolio, account, market observations,
and lifecycle state; \(A_t\) contains resolved actions; \(M_t\) records which
actions are presently admissible; and \(p_t\) identifies the definitions,
observations, assumptions, and prior journal state.  The policy consumes this
object but cannot edit its authority rules.

Today these pieces are separate.  \texttt{PortfolioState} contains exact
position lifecycle, cash, underlying units, sequence, and observation time;
\texttt{CandidateGrammar} compiles bounded candidate actions; and
\texttt{PortfolioGraphSample} provides a static portfolio graph.  There is no
combined account-state graph or legal-mask artifact in the public release.

A proposed environment transition binds an action to versioned state:
\[
  (s_{t+1},r_t,e_t)
    =\mathcal T_{\omega}(s_t,a_t,\xi_{t+1}),
\]
where \(\mathcal T_{\omega}\) combines the Volt runtime with declared market
and execution model \(\omega\), \(\xi_{t+1}\) contains the next market and
external events, \(r_t\) is a derived reward record, and \(e_t\) is immutable
transition evidence.  The implemented runtime itself accepts a
\texttt{PortfolioState}, \texttt{TransitionAction}, and
\texttt{MarketObservation}, returning the next state and one journal entry.
Replaying the same versioned inputs reproduces state and emissions.

\section{Actions and masks}

The initial grammar contains \texttt{AddLeg}, \texttt{RemoveLeg},
\texttt{ResizeLeg}, \texttt{Wait}, and holder-authorized \texttt{Exercise}.
The runtime action vocabulary is distinct:
\texttt{Hold}, \texttt{Close}, \texttt{Exercise}, \texttt{Assign}, and
\texttt{Expire}.  No current compiler turns close, roll, hedge, take-profit,
or reduce-risk instructions into a multi-step runtime plan.
Assignment remains an external event because a short writer does not possess
the holder's exercise right.

Action resolution takes place at the decision clock.  A semantic request such
as ``roll the short call to the next monthly expiration near 25 delta'' must
resolve only against contracts and observations visible at that time.  The
result names exact listed contracts and quantities.  Volt enforces structural
and exercise rules; an application must calculate estimated resource use and
apply liquidity, position, and account constraints before a policy normalizes
its scores.

Masks may change abruptly along a path.  Expiration removes a contract;
exercise consumes a right; a fill changes available cash; assignment creates
inventory; a market halt may make an otherwise legal close unavailable under
the execution model.  These changes are part of the environment rather than
exceptions handled in model code.

\section{Rewards and horizons}

A control environment may derive financial outcomes as a vector:
\[
  r_t=(\Delta\text{equity},\text{cost},\text{drawdown},
       \text{capital},\text{turnover},\text{constraint events}).
\]
The current simulator records fills, costs, journals, marked pathwise P\&L,
leg attribution, and terminal P\&L.  Drawdown, capital, turnover, and
constraint-event components require additional named models.  An experiment
may scalarize its declared vector for policy optimization.  A common form is
\[
  R_t=\Delta\text{equity}_t
      -\lambda_c\text{cost}_t
      -\lambda_d\Delta\text{drawdown}_t
      -\lambda_k\text{capital}_t,
\]
with coefficients fixed in the manifest.  Reward shaping requires particular
care because a policy can optimize the shaping rule while degrading terminal
wealth or tail exposure.

Economic and computational horizons should be named separately.  The episode
may end at expiration, liquidation, a risk breach, or the end of the historical
window.  Training can still use shorter \(n\)-step returns or truncated
recurrent unrolls with a bootstrap value:
\[
  \widehat R_t^{(n)}
    =\sum_{j=0}^{n-1}\gamma^j R_{t+j}
     +\gamma^n\widehat V(s_{t+n}).
\]
This truncation limits computation and credit-assignment length; it does not
declare the financial position finished at \(t+n\).  The bootstrap state and
all open obligations remain in the trajectory record.

\part{Learning and Search}

\section{Policy architecture}

A heterogeneous graph encoder maps \(G_t\) to a portfolio state \(h_t\).
Candidate actions receive embeddings \(e(a)\), and a policy head scores each
pair:
\[
  \ell_\theta(s_t,a)
    =f_\theta\!\left(h_t,e(a),h_t\odot e(a)\right).
\]
Masked logits define
\[
  \pi_\theta(a\mid s_t)
    =\operatorname{softmax}_{a\in A_t}
      \bigl(\ell_\theta(s_t,a)+\log M_t(a)\bigr),
\]
using the implementation convention that excluded actions receive negative
infinity before normalization.  A value head estimates expected discounted
return or its distribution.  The encoder may be initialized from supervised
strategy ranking, although control performance remains the criterion that
determines whether the transfer was useful.

Policy-gradient, actor--critic, or value-based algorithms can consume this
interface.  The initial implementation should favor a well-understood
baseline such as PPO with generalized advantage estimation
\cite{gae,ppo}; algorithmic novelty would make failures harder to locate
while the environment and reward contract are still being tested.

\section{Trajectory collection}

\begin{center}
\begin{tikzpicture}[node distance=7mm]
  \node[controlbox] (state) {\textbf{Proposed state}\\graph, candidates, mask,
    evidence};
  \node[controllearn,right=of state] (policy) {\textbf{Policy}\\score and
    sample an admissible action};
  \node[controlbox,right=of policy] (runtime) {\textbf{Runtime}\\fills, costs,
    events, exercise, settlement};
  \node[controlbox,below=8mm of runtime] (store) {\textbf{Trajectory
    store}\\states, actions, rewards, propensities};
  \node[controllearn,left=of store] (learner) {\textbf{Learner}\\policy, value,
    and diagnostics};
  \draw[controlflow] (state) -- (policy);
  \draw[controlflow] (policy) -- (runtime);
  \draw[controlflow] (runtime) -- (store);
  \draw[controlflow] (store) -- (learner);
  \draw[controlfeedback] (runtime.north) -- ++(0,5mm) -|
    node[pos=0.72,above,font=\sffamily\scriptsize,text=SNMuted]
    {next state} (state.north);
  \draw[controlfeedback] (learner.west) -- ++(-6mm,0) |-
    node[pos=0.25,left,font=\sffamily\scriptsize,text=SNMuted]
    {checkpoint} (policy.south);
\end{tikzpicture}
\begingroup
\captionsetup{hypcap=false}
\captionof{figure}{The proposed control boundary.  Policy workers would choose
from compiler masks; the runtime owns portfolio transitions and the learner
would update from versioned trajectories.}
\label{fig:control-loop}
\endgroup
\end{center}
\par\medskip

Independent workers can collect trajectories for different dates,
underlyings, scenario paths, execution assumptions, and random seeds.  They
share a frozen policy checkpoint during a collection epoch and emit immutable
trajectory fragments.  A deterministic merge orders fragments by stable
identifiers before optimization, so worker completion order does not change
the training batch.

The stored record includes the policy version and action probability used at
collection time.  Those fields matter when asynchronous workers lag the
learner or when an allocation rule changes which experiences enter an update.
Without them, off-policy drift and selection bias become difficult to measure
or correct.

\section{Rollout allocation}

\subsection{Search frontier}

Game-playing systems demonstrate that a learned policy and value function can
concentrate tree-search effort without simulating every continuation to a
terminal state \cite{alphazero,muzero}.  An option-management tree has related
structure but a different source of branching: market scenarios, external
events, execution uncertainty, and portfolio actions all create descendants.
The cost of one expansion can also vary sharply with the scenario model and
remaining horizon.

An allocator receives a frontier of partial paths and a finite budget.  It may
decide which nodes receive another action expansion, another market scenario,
or a longer continuation.  Useful signals include policy entropy, value
uncertainty, tail-risk uncertainty, novelty, estimated improvement, path
probability, and simulation cost.  A simple upper-confidence or tree-search
rule should establish the baseline before a combinatorial allocator is added.

\subsection{QUBO proposal}

Let \(x_i\in\{0,1\}\) indicate whether frontier item \(i\) receives additional
compute.  For equal-cost expansions and cardinality budget \(k\), consider
\[
  \min_x
    -\sum_i u_i x_i
    +\lambda\sum_{i<j}s_{ij}x_ix_j
    +\rho\left(\sum_i x_i-k\right)^2.
\]
The linear utility \(u_i\) can combine expected improvement, epistemic
uncertainty, tail relevance, and path probability.  Pairwise similarity
\(s_{ij}\) penalizes allocating the whole batch to nearly identical portfolio
states or market paths.  The final term encodes the compute budget.  Variable
simulation costs require a weighted constraint and binary slack variables, or
a separate constrained quadratic formulation.

This construction has a close published precedent.  Kholodov, Salloum, Jnadi,
Khubiev, and Petrenko formulate Monte Carlo episode selection as QUBO, using
linear terms for high-return episodes and quadratic penalties for trajectory
redundancy and poor state coverage \cite{quboepisodes}.  Their evidence comes
from a finite GridWorld and establishes feasibility of the selection
mechanism.  Portfolio control requires its own evidence.

QUBO can enter on either side of simulation.  Frontier allocation occurs
before simulation and can save runtime cost, but its utility estimates are
necessarily incomplete.  Replay selection occurs after collection, when
returns and coverage are known; it saves learner compute while leaving rollout
cost unchanged.  Prioritized experience replay provides the natural
non-QUBO comparison for the latter \cite{per}.

\subsection{Coverage and bias}

Aggressive selection can remove the paths that reveal delayed losses, rare
assignment states, or a mistaken value estimate.  A production-quality
experiment therefore reserves part of every budget for random or
stratified exploration and treats rejection as a decision for the current
budget, never as proof that a legal path is irrelevant.

Every selected item records its score components, pairwise penalties,
selection probability where estimable, solver configuration, and QUBO
solution quality.  Replay-based learning can use importance weights to reduce
the bias introduced by nonuniform sampling.  Frontier selection is harder:
unexpanded descendants have no observed return, so evaluation must compare
the resulting policy with exhaustive small trees and with fixed-budget
baselines.

The quadratic similarity matrix grows as \(O(n^2)\).  Large frontiers will
need clustering, sparse neighborhoods, or staged selection before QUBO
construction.  Solver time is charged to the same wall-clock budget as
simulation; a sophisticated allocator that costs more than the rollouts it
saves has failed its engineering objective.

\section{Parallel execution}

Scenario paths and environment instances form the main parallel dimension.
Graph encoding and policy inference can be batched on a GPU, while homogeneous
pricing or scenario kernels may share that device.  Exercise, assignment,
correction, and journal processing remain control-heavy and can stay on the
CPU until profiling supports a different lowering.

Rollout scheduling may be hierarchical.  A coordinator divides a global
budget among market slices or scenario families; local allocators choose
frontier nodes within each partition; workers execute the selected requests.
Stable identifiers preserve the relationship between a partial path, its
parent, the policy checkpoint, and the evidence returned by the runtime.

\part{Evaluation}

\section{Protocol}

Development begins with synthetic and bounded historical environments in
which exhaustive or near-exhaustive rollout trees remain available.  The
protocol then expands along one dimension at a time: horizon, action grammar,
scenario count, execution uncertainty, and number of contracts.  This order
makes it possible to attribute a performance change to a specific source of
branching.

Train, validation, and test periods remain chronological.  Reward
coefficients, termination rules, truncation length, bootstrap convention,
random exploration share, and rollout budget are frozen before test
evaluation.  Results are reported over multiple policy and environment seeds,
including failures and constraint events.

\section{Baselines}

The control policy is compared with economically meaningful management rules:
hold to expiration, fixed take-profit and stop rules, clock-based hedging,
delta-threshold hedging, and simple roll conventions.  Learning comparisons
include uniform rollout allocation, random frontier expansion, beam search,
UCT-style tree search \cite{uct}, greedy utility, greedy utility with diversity,
and prioritized replay.  QUBO receives the same candidate frontier and total
compute allowance.

Measurements cover:
\begin{itemize}[nosep]
  \item net and gross P\&L, drawdown, tail loss, turnover, and capital use;
  \item frequency and severity of constraint or lifecycle failures;
  \item policy stability across costs, spreads, and scenario assumptions;
  \item state and action coverage, including assignment and exercise paths;
  \item value calibration and bootstrap error by remaining horizon;
  \item simulator calls, expanded nodes, learner updates, and wall-clock time;
  \item allocation overhead and solution quality for each scheduler.
\end{itemize}

Equal simulator counts answer a statistical-efficiency question; equal
wall-clock budgets answer an engineering question.  Both are required because
parallel collection, GPU batching, and QUBO solution overhead change the
relationship between those budgets.

\section{Evidence and boundary}

The policy should choose only actions authorized by the future compiler and
application mask.  The scenario layer supplies declared close fills and fees;
the runtime validates exact contract actions, handles exercise and assignment,
updates cash and inventory, and journals the transition.  The
learner owns scores, action probabilities, value estimates, advantages, and
checkpoint updates.  Rollout allocation owns compute scheduling.  These
records meet through stable state and request identifiers without sharing
authority.

The repository contains no quote-backed historical policy result.  A touch
fill remains a declared execution convention, and a path selected after
observing its return cannot be described as a tradable policy.  Economic
evidence requires an untouched period, capacity and cost analysis, sensitivity
to market and assignment models, and comparison with simple management rules.

\section{Conclusion}

Volt supplies the deterministic state and transition core from which a
sequential-control environment can be built while contract semantics remain
stable.  That boundary can let research move from PPO to tree search or
QUBO-guided scheduling without changing what an option, exercise, fill, or
assignment means.  The immediate research goal is a small replayable
environment with exhaustive controls; only after its learning and accounting
agree should rollout allocation be asked to save compute on larger trees.

\begin{thebibliography}{9}

\bibitem{gae}
J. Schulman, P. Moritz, S. Levine, M. Jordan, and P. Abbeel,
\emph{High-Dimensional Continuous Control Using Generalized Advantage
Estimation},
International Conference on Learning Representations, 2016,
\url{https://arxiv.org/abs/1506.02438}.

\bibitem{ppo}
J. Schulman, F. Wolski, P. Dhariwal, A. Radford, and O. Klimov,
\emph{Proximal Policy Optimization Algorithms},
2017,
\url{https://arxiv.org/abs/1707.06347}.

\bibitem{alphazero}
D. Silver et al.,
\emph{Mastering Chess and Shogi by Self-Play with a General Reinforcement
Learning Algorithm},
2017,
\url{https://arxiv.org/abs/1712.01815}.

\bibitem{muzero}
J. Schrittwieser et al.,
\emph{Mastering Atari, Go, Chess and Shogi by Planning with a Learned Model},
Nature, 588, pp.~604--609, 2020,
\url{https://doi.org/10.1038/s41586-020-03051-4}.

\bibitem{per}
T. Schaul, J. Quan, I. Antonoglou, and D. Silver,
\emph{Prioritized Experience Replay},
International Conference on Learning Representations, 2016,
\url{https://arxiv.org/abs/1511.05952}.

\bibitem{uct}
L. Kocsis and C. Szepesvári,
\emph{Bandit Based Monte-Carlo Planning},
European Conference on Machine Learning, 2006,
\url{https://doi.org/10.1007/11871842_29}.

\bibitem{quboepisodes}
Y. A. Kholodov, H. Salloum, A. Jnadi, K. Yu. Khubiev, and A. Petrenko,
\emph{Quantum-Inspired Episode Selection for Monte Carlo Reinforcement
Learning via QUBO Optimization},
Computer Research and Modeling, 18(2), pp.~273--288, 2026,
\url{https://doi.org/10.20537/2076-7633-2026-18-2-273-288}.

\end{thebibliography}

\end{document}
